\documentclass[10pt,a4paper]{amsart}
\usepackage[margin=19mm]{geometry}
\usepackage{amsmath,amssymb,mathtools}
\usepackage[scr=rsfs,cal=euler]{mathalfa}
\usepackage{xfrac}
\usepackage[unicode,hidelinks,bookmarks=false]{hyperref}
\newcommand{\ie}{i.e.\ }
\newcommand{\cf}{cf.\ }
\newcommand{\resp}{respectively\ }
\newcommand{\Subsection}{\subsection}
\providecommand{\nobreakdash}{\nobreak\hbox{-}}
\setlength{\emergencystretch}{2em}
\multlinegap0pt
\usepackage{amssymb}
\usepackage{xcolor}
\usepackage{algorithm}\usepackage{algpseudocode}
\newtheorem{lemma}{Lemma}
\title{Quantitative Estimates for Mean-Field Limits and Correlation Functions through a Duality Framework}
\author{Nadia Khoury, Pierre-Emmanuel Jabin}
\date{Source: arXiv:2605.02058v2 (CC BY 4.0)}
\begin{document}
\maketitle

\noindent\emph{Source and licence.} Quantitative Estimates for Mean-Field Limits and Correlation Functions through a Duality Framework. Version arXiv:2605.02058v2 (CC BY 4.0), distributed by arXiv under the Creative Commons Attribution 4.0 International licence, http://creativecommons.org/licenses/by/4.0/, as recorded in the arXiv licence metadata for this exact version and shown on the abstract page https://arxiv.org/abs/2605.02058v2. Full licence notice: \texttt{licenses/arxiv-2605.02058-CC-BY-4.0.txt}.\par\smallskip

\noindent\textit{Editorial scope.} Counted unit: the article's lemma lem:one\_particle\_k (source0.tex lines 601--615). The article's complete original proof of it is delivered with it (source0.tex lines 617--794, one \texttt{proof} environment of 178 lines). No mathematics is written, repaired or re-derived: the statement and the proof are the article's own lines, in the article's own notation, and no retained line is substituted or re-flowed.  No further source-local context is needed: every cross-reference of every retained line resolves inside the two delivered spans.  Nothing was omitted from any retained span, so the package carries no omission record.  No retained line contains a citation, so the wrapper carries no bibliography and no bibliography entry.  No retained line contains a figure environment, an image-inclusion command or drawing code, so no image is delivered and the package declares no asset dependency.  Editorial formatting only, in addition: an AMS \texttt{amsart} preamble that restates the article's own declarations and macros used by the retained lines, namely the environments \texttt{lemma} on separate counters, together with the standard packages those lines need.  The retained environments print in this document as Lemma 1 in order of appearance, whereas the article prints the same items with the numbering of its own full document.  The complete native source of the article as distributed in the arXiv e-print for this exact version is bundled unchanged as \texttt{sources/199688/source0.tex}.
\begin{lemma} \label{lem:one_particle_k} 
Fix $i\in\{1,\dots,n\}$ and $k\ge 1$.
Let
\[
A_n^{(i)} := v_i\cdot\nabla_{x_i} + (K*f)(x_i)\cdot\nabla_{v_i},
\]
with $K \ast f \in L^{\infty}\left(0,T ;W^{k+1,\infty}(\Omega ; \mathbb{R}^d) \right)$, and denote by $G^{(n,i)}_{t,T}$ the corresponding operator solving the PDE: $( \partial_t+A_n^{(i)} )u=0$.
Then, there exists a family of operators $G^{(n,i)}$ and $\{G^{(n,i),\ell,k}\}_{\ell=0}^k$
bounded over $L^2$ with an operator norm that depends only on $k$ and $T-t$
such that 
\[
G_{t,T}^{(n,i)}\nabla^k_{z_i} g_{n,i}=\sum_{\ell =0}^k\nabla_{z_i}^{\,\ell}
G^{(n,i),\ell,k}_{t,T} g_{n,i}.
\]
\end{lemma}

\begin{proof}
Note that all variables $x_j$ and $v_j$ for $j\neq i$ are just parameters so that we will not indicate them in the various formulas. Using the method of characteristics, we define the characteristic curves
\[
\dot{X}_i(t,s,x,v) = V_i(t,s,x,v),
\qquad
X_i(t=s,s,x,v)=x_i,
\]
and
\[
\dot{V}_i(t,s,x,v) = K * f \bigl(X_i(t,s,x,v)\bigr),
\qquad
V_i(t=s,s,x,v) = v_i.
\]
Define
\[
F_{n,i}(t) := f_{n,i}\bigl(t, X_i(t,s,x,v), V_i(t,s,x,v)\bigr).
\]
By the chain rule, we compute
\begin{align*}
\partial_t F_{n,i}
&=
\partial_t f_{n,i}(t,X_i,V_i)
+
\bigl(\dot X_i\cdot\nabla_{x_i}
+
\dot V_i\cdot\nabla_{v_i}\bigr)
f_{n,i}(t,X_i,V_i) \\
&=
\partial_t f_{n,i}(t,X_i,V_i)
+
A_n^{(i)} f_{n,i}(t,X_i,V_i) \\
&=
\nabla_{z_i}^k g_{n,i}(t,X_i,V_i).
\end{align*}
Integrating from $t$ to $T$ yields
\[
F_{n,i}(t)
=
F_{n,i}(T)
-
\int_t^T
\nabla_{z_i}^k g_{n,i}\bigl(\tau, X_i(\tau,s,x,v), V_i(\tau,s,x,v)\bigr)\,d\tau.
\]
Taking $s=t$, we obtain the representation of the operator $G^{(n,i)}_{t,T}$ as
\[
G^{(n,i)}_{t,T} \bar f_{n,i}=\bar f_{n,i}(X_i(T,t,x,v),V_i(T,t,x,v)),
\]
so that $G^{(n,i)}_{t,\tau}$ is naturally bounded on $L^2$.

\noindent
We now need to rewrite $G^{(n,i)}_{t,\tau}\nabla_{z_i}^k$ as a finite sum of
derivatives acting outside the operator, and we proceed by induction on $k$.

\medskip
\noindent\textbf{Step 1: Case $k=1$.}
Assume that $g_{n,i}$ is sufficiently smooth and denote by
\[
Z_i(\tau,t,x,v) := (X_i(\tau,t,x,v), V_i(\tau,t,x,v))
\]
the characteristic flow associated with the one--particle operator $A_n^{(i)}$.
By the definition of the backward flow operator, for any test function $h$ we have
\[
\big(G^{(n,i)}_{t,\tau} h\big)(x,v) = h\big(\tau, Z_i(\tau,t,x,v)\big).
\]
In particular,
\[
G^{(n,i)}_{t,\tau}\nabla_{z_i} g_{n,i}
=
\big(\nabla_{z_i} g_{n,i}\big)\big(\tau,Z_i(\tau,t,x,v)\big).
\]
The key observation is that, since the right--hand side of the equation already
contains one derivative, we do not aim to eliminate derivatives altogether, but
rather to rewrite this expression so that exactly one derivative appears
\emph{outside} an operator acting on $g_{n,i}$. By the chain rule,
\[
\nabla_{z_i}\!\left(g_{n,i}\big(\tau,Z_i\big)\right)
=
\big(\nabla_{z_i} g_{n,i}\big)\big(\tau,Z_i\big)\,\nabla Z_i,
\]
where $\nabla Z_i$ denotes the Jacobian matrix of the flow with respect to $z_i$.
Since the characteristic flow is a diffeomorphism, the matrix $\nabla Z_i$ is
invertible, and we may therefore write
\[
\big(\nabla_{z_i} g_{n,i}\big)\big(\tau,Z_i\big)
=
(\nabla Z_i)^{-1}
\nabla_{z_i}\!\left(g_{n,i}\big(\tau,Z_i\big)\right).
\]
Applying the product rule to the right--hand side yields the decomposition
\[
\big(\nabla_{z_i} g_{n,i}\big)\big(\tau,Z_i\big)
=
\nabla_{z_i}\!\left((\nabla Z_i)^{-1} g_{n,i}\big(\tau,Z_i\big)\right)
-
\nabla_{z_i}(\nabla Z_i)^{-1}\,
g_{n,i}\big(\tau,Z_i\big).
\]
This naturally leads us to define the operators
\[
G^{(n,i),1,1}_{t,\tau} g_{n,i}
:=
(\nabla Z_i)^{-1}\, g_{n,i}\big(\tau,Z_i\big),
\qquad
G^{(n,i),0,1}_{t,\tau} g_{n,i}
:=
-\nabla_{z_i}(\nabla Z_i)^{-1}\, g_{n,i}\big(\tau,Z_i\big),
\]
so that
\begin{equation}\label{eq2}
G^{(n,i)}_{t,\tau}\nabla_{z_i} g_{n,i}
=
\nabla_{z_i}\big(G^{(n,i),1,1}_{t,\tau} g_{n,i}\big)
+
G^{(n,i),0,1}_{t,\tau} g_{n,i}.
\end{equation}
The matrices $(\nabla Z_i)^{-1}$ and $\nabla_{z_i}(\nabla Z_i)^{-1}$ depend only on
the characteristic flow. Under the assumption that $K*f \in L^{\infty}(0,T;W^{2,\infty}(\Omega;\mathbb{R}^d))$, standard Gr\"onwall estimates for the
Jacobian of the flow imply that the coefficients $(\nabla Z_i)^{-1}$ and
$\nabla_{z_i}(\nabla Z_i)^{-1}$ are bounded uniformly in space on any finite time
interval. Consequently, for $\ell=0,1$, all $0\le t\le \tau\le T$, and all $h \in L^2$, the operators
$G^{(n,i),\ell,1}_{t,\tau}$ are bounded on $L^2$, with
\[
\|G^{(n,i),\ell,1}_{t,\tau} h\|_{L^2}
\le C_1\,e^{C_1\,(T-t)}\|h \|_{L^2}.
\]
Inserting identity \eqref{eq2} into the solution representation with $k=1$
yields the desired result.

\medskip
\noindent\textbf{Step 2: General case.}
Let $G^{(n,i)}_{t,\tau}\nabla_{z_i}^{k} g_{n,i}$ for a sufficiently smooth $g_{n,i}$. We prove the corresponding statement at order $k$.
%Consider $G^{(n,i)}_{t,\tau}\nabla_{z_i}^{k} g_{n,i}$ and write
%\[
%G^{(n,i)}_{t,\tau}\nabla_{z_i}^{k} g_{n,i}
%=
%G^{(n,i)}_{t,\tau}\nabla_{z_i}\big(\nabla_{z_i}^{k-1} g_{n,i}\big).
%\]
Applying the $k=1$ decomposition of Step~1 to the function
$h:=\nabla_{z_i}^{k-1}g_{n,i}$ gives
\begin{equation}\label{eq:step1_applied_to_h}
G^{(n,i)}_{t,\tau}\nabla_{z_i}h
=
\nabla_{z_i}\big(G^{(n,i),1,1}_{t,\tau}h\big)
+
G^{(n,i),0,1}_{t,\tau}h,
\end{equation}
hence
\[
G^{(n,i)}_{t,\tau}\nabla_{z_i}^{k} g_{n,i}
=
\nabla_{z_i}\Big(G^{(n,i),1,1}_{t,\tau}\nabla_{z_i}^{k-1}g_{n,i}\Big)
+
G^{(n,i),0,1}_{t,\tau}\nabla_{z_i}^{k-1}g_{n,i}.
\]
We now repeatedly apply the first-order decomposition from Step~1 to move
derivatives outside the flow operators. Since each application produces only
bounded coefficient operators coming from the characteristic flow and lowers the
number of derivatives falling on $g_{n,i}$ by one, after finitely many iterations
we may expand and re-index the resulting finite sum to obtain
\[
G^{(n,i)}_{t,\tau}\nabla_{z_i}^{k} g_{n,i}
=
\sum_{\ell=0}^{k}
\nabla_{z_i}^{\,\ell}\big( G^{(n,i),\ell,k}_{t,\tau} g_{n,i}\big),
\]
for a suitable family of bounded operators
$\{G^{(n,i),\ell,k}_{t,\tau}\}_{\ell=0}^k$.

\noindent
Finally, by the same reasoning as in Step 1, for $0 \leq \ell \leq k$, each $G^{(n,i),\ell,k}_{t,\tau}$ is bounded
on $L^2$, and we may write
\[
\|G^{(n,i),\ell,k}_{t,\tau} h\|_{L^2}
\le C_k(T-t)\,\|h\|_{L^2},
\qquad \forall\, h\in L^2.
\]
This completes the induction step.
\end{proof}

\end{document}
